\uuid{x7jq}
\exo7id{518}
\titre{exo7 518}
\auteur{ridde}
\organisation{exo7}
\datecreate{1999-11-01}
\isIndication{false}
\isCorrection{true}
\chapitre{Suite}
\sousChapitre{Convergence}
\cpge{mpsi-mp2i.nombres-reels-suites.c, pcsi.nombres-reels-suites.b, ptsi.nombres-reels-suites.b, mpsi-mp2i.nombres-reels-suites.d, pcsi.nombres-reels-suites.c, ptsi.nombres-reels-suites.c, mpsi-mp2i.nombres-reels-suites.f, pcsi.nombres-reels-suites.e, ptsi.nombres-reels-suites.e, bcpst1.suites-reelles}
\module{Analyse}
\niveau{L1}
\difficulte{}

\contenu{
\question{
\'Etudier la convergence des suites :\\

$\sqrt{n^2 + n + 1}-\sqrt n \qquad \dfrac{n\sin (n)}{n^2 + 1} \qquad \dfrac1n  + (-1)^n
\qquad n \sum\limits_{k = 1}^{2n + 1}\dfrac 1{n^2 + k} \qquad
\dfrac 1n \sum\limits_{k = 0}^{n-1}\cos (\dfrac1{\sqrt{n + k}})$
}
\reponse{
Suite non convergente car non bornée.
Suite convergente vers $0$.
Suite non convergente car la sous-suite $u_{2p} = 1+\frac {1}{2p}$ est
toujours plus grande que $1$. Alors que la sous-suite $u_{2p+1} =
-1+\frac {1}{2p+1}$ est toujours plus petite que $0$.
}
}
